\documentclass[../../../include/open-logic-section]{subfiles}

\begin{document}

\olfileid[fr]{sth}{story}{blv}

\olsection{Annexe : la loi fondamentale V de Frege}

Dans la \olref[rus]{sec}, nous avons expliqué que Russell avait
formulé son paradoxe comme une difficulté pour le système exposé par
Frege dans ses \emph{Grundgesetze}. Le système de Frege ne comprenait
pas de formulation directe de la compréhension naïve. Dans cette
annexe, nous expliquerons donc très brièvement ce qu'il
\emph{comprenait}, puis son rapport avec la compréhension naïve et
avec le paradoxe de Russell.

Le système de Frege relève de la \emph{logique du second ordre} et
devait formaliser la notion d'\emph{extension d'un concept}.\footnote{À
strictement parler, Frege cherche à formaliser une notion plus
générale : le «~parcours de valeurs~» d'une fonction. Les extensions
de concepts en constituent un cas particulier. Pour les détails,
voir \citet[pp.\ 8--9]{Heck2012}.} En nous inspirant de la notation de
Frege, nous écrirons $\fregeext{x}{F(x)}$ pour \emph{l'extension du
concept~$F$}. Ce dispositif prend un \emph{prédicat}, «~$F$~», et
en fait un \emph{terme} (du premier ordre),
«~$\fregeext{x}{F(x)}$~». Frege s'en est servi pour proposer la
\emph{définition} suivante de l'appartenance :
\[
	a \in b =_\text{df} \exists G(b = \fregeext{x}{G(x)} \land Ga)
\]
En gros, $a \in b$ si et seulement si $a$ tombe sous un concept dont
l'extension est $b$. (Le quantificateur «~$\exists G$~» est du second
ordre.) Frege soutenait également le principe suivant, connu sous le
nom de \emph{loi fondamentale V} :
$$\fregeext{x}{F(x)} = \fregeext{x}{G(x)} \liff \forall x (Fx \liff Gx)$$
En gros, deux concepts ont la même extension si et seulement s'ils
sont coextensifs. (Ici encore, «~$F$~» et «~$G$~» sont en position de
prédicat.) Un principe simple relie maintenant l'appartenance à la
satisfaction d'une propriété :

\begin{lem}[dans les \emph{Grundgesetze}]\ollabel{lem:Fregeextensions}
$\forall F \forall a(a \in \fregeext{x}{F(x)} \liff Fa)$
\end{lem}

\begin{proof}
Fixons $F$ et $a$. Alors $a \in \fregeext{x}{F(x)}$ si et seulement
si $\exists G(\fregeext{x}{F(x)}
= \fregeext{x}{G(x)} \land Ga)$ (par définition de l'appartenance),
si et seulement si $\exists G(\forall x(Fx \liff Gx) \land Ga)$
(par la loi fondamentale V), si et seulement si $Fa$ (par la logique
élémentaire du second ordre).
\end{proof}

On en déduit presque aussitôt la compréhension naïve :

\begin{lem}[dans les \emph{Grundgesetze}.]
$\forall F \exists s \forall a (a \in s \liff Fa)$
\end{lem}

\begin{proof}
Fixons $F$. Le \olref{lem:Fregeextensions} donne alors $\forall a (a \in
\fregeext{x}{F(x)} \liff Fa)$ ; on obtient donc
$\exists s\forall a(a \in s \liff Fa)$ par généralisation
existentielle. Puisque $F$ était arbitraire, le résultat s'ensuit.
\end{proof}

Le paradoxe de Russell s'ensuit si l'on prend $F$ tel que
$\forall x(Fx \liff x \notin x)$.

\end{document}
